\chapter{Limitations and Future Work}
\label{chap:future}

\section{Current Limitations}

\subsection{2D Processing and Inter-Slice Context}

The most fundamental limitation of this work is the slice-by-slice processing paradigm. Each
axial slice is segmented independently, which means the model has no access to the anatomical
context provided by neighbouring slices. This causes the boundary prediction to degrade at the
apex and base of curved organs --- the high HD95 for spleen (25.7~mm) and right kidney
(26.1~mm) are the clearest manifestation of this problem. A 3D model or a 2.5D approach
(processing 2D slices with an additional context channel from adjacent slices) would address
this directly.

\subsection{Dataset Scale}

Fifty CT volumes is a modest dataset for a 13-class segmentation problem. The ablation study
shows a clear performance gradient with training set size: the jump from 30 to 50 volumes adds
more than $+$0.07 Dice. It is reasonable to expect further gains with 100 or 200 volumes, which
are now available through FLARE23 and other public benchmarks. Transfer learning from a
self-supervised pretrained backbone (e.g., a ResNet or ViT pretrained on large-scale CT
datasets) could partially offset the data scarcity issue.

\subsection{Fixed Input Resolution}

All CT slices are resized to $256 \times 256$. For some high-resolution volumes ($> 512 \times
512$ native resolution), this downsampling discards fine boundary detail. Adaptive tiling
--- processing each volume at its native resolution in non-overlapping patches and stitching
predictions --- would recover this detail without the memory cost of upsampling the entire
model.

\subsection{No Post-Processing}

The current pipeline outputs raw argmax predictions without any spatial post-processing (e.g.,
connected-component analysis to remove spurious small predictions, or conditional random fields
for boundary refinement). For clinical use, such post-processing steps are standard and would
improve HD95 scores on organs like the esophagus and pancreas, which occasionally produce
disconnected prediction islands.

\section{Future Directions}

\subsection{2.5D and Pseudo-3D Architectures}

The most immediate extension is to move from pure 2D to a 2.5D input --- stacking the target
slice with its immediate axial neighbours to give the model three-slice context. This approach
is memory-efficient and has been shown to recover most of the performance gap between 2D and
full 3D models on abdominal CT data.

\subsection{Semi-Supervised Learning}

FLARE22 includes an additional 2{,}000 unlabelled CT volumes. A semi-supervised approach ---
using the labelled 50 volumes for supervised training and the unlabelled 2{,}000 for consistency
regularisation or pseudo-labelling --- could dramatically expand the effective training set
without the cost of manual annotation.

\subsection{Uncertainty Estimation}

The model currently produces a hard argmax prediction with no indication of confidence. Adding
Monte Carlo dropout at inference (leaving the bottleneck Dropout2d active during evaluation)
would produce per-voxel uncertainty estimates at negligible cost. Uncertainty maps are
clinically useful --- they highlight regions where the model is unsure and a radiologist should
verify the segmentation.

\subsection{Transformer-Based Backbones}

Recent architectures such as Swin-UNETR replace convolutional encoders with Swin Transformers,
capturing long-range spatial dependencies that convolutions miss. While these models require
more data to train from scratch, fine-tuning a publicly available Swin-UNETR checkpoint on
FLARE22 data would be feasible within the same hardware constraints as this project.

\subsection{Multi-Dataset Pretraining}

Training a shared backbone on several abdominal CT datasets (e.g., KiTS23, MSD-Liver,
CHAOS) before fine-tuning on FLARE22 could improve robustness to scanner variability and
organ co-occurrence patterns. This multi-dataset strategy is increasingly standard in
large-scale medical imaging benchmarks and represents a natural next step for this work.
